\chapter*{写在前面}
\addcontentsline{toc}{chapter}{写在前面}

自由电子经过光场附近时，为什么会出现以光子能量为间隔的边带？如果场具有空间旋涡、多色相位或量子涨落，电子究竟带走了什么信息？测量到的能谱又能在多大程度上反推出电子与材料的相互作用？这些问题必须沿着同一条物理链回答：先从谱峰位置和强度看见电子怎样与近场交换能量；随后沿轨迹求出周期相位、完成束流和探测器的前向计算，并让相干边带在自由漂移中形成阿秒结构。对这些结果的近似范围有了具体判断之后，再回到 Dirac 方程和 Maxwell 散射，求出实际材料场以及模型误差，最后讨论量子光、开放系统与多电子关联。

本书从 PINEM 实验中可直接看到的能量增益与损失边带出发，先建立读图和估计耦合强度所需的最小模型。第一部分从轨迹相位、谱矩、束流平均和自由传播完成一次从近场到可观测量的计算；第二部分再检验电子窄带与无反冲近似，并通过 Maxwell 散射、球形 Mie 场和有限脉冲研究实际结构的输入场。能量--角动量合成维度、量子光统计、退相干、电子交换和测量反演在这些对象已经具备实际意义之后依次出现。电子测量还面临另一个贯穿多章的问题：即使把已有数据测得更精确，不同近场也可能产生同一组可观测结果。理解这种不可辨识性，再设计新的 OAM 通道、参考场或交换历史标记，是从 PINEM 图像走向量子信息测量的自然延伸。章节中的计算图和章末复现说明使用明确的数学模型。现象学扫描与真实实验的绝对标定属于不同证据层级；有关参数应按具体模型的适用条件解读。

全书采用 SI 单位制，$e>0$ 是元电荷大小，电子电荷 $q_e=-e$。实场的频域约定为 $E(t)=\operatorname{Re}[E_\omega\ee^{-\ii\omega t}]$，正频部分因此为 $E^{(+)}=E_\omega\ee^{-\ii\omega t}/2$。采用这一约定时，经典 PINEM 复耦合的系数是 $e/(2\hbar\omega)$；若另一本参考文献把正频分量本身称作“复场振幅”，其表面上缺少的 $1/2$ 因子只是幅度记号的差异，不是新的物理效应。后文均按此约定计算数值振幅和边带相位。
